\section{Results}

\subsection{Evaluation Metrics}
Multiple metrics are used to evaluate the model performance, here in this section we outline each metric and give a motivation as to why. Throughout the results and evaluation both metrics are
reported together. The reason as to why both metrics are reported is that even though the geodesic angle error might decrease the yaw error may increase.

\subsubsection{Geodesic Angle Error}
To determine the error between the two quaternions that exist as rotations on SO(3), we use the geodesic rotation angle between the two
quaternions. This metric measures the minimal rotation needed to align the two quaternions.

\begin{equation}
    \epsilon_{\text{geo}} = 2\arccos\!\bigl(|\hat{q} \cdot q_{\text{gt}}|\bigr)
\end{equation}

The reason as to why we use the geodesic angle error is that it is the correct distance between two quaternions on SO(3). A Euclidean difference is not rotation invariant meaning 
$q$ and $-q$ would represent different rotations, where they are the same rotation on SO(3). It also captures the total orientation error (roll, pitch, and yaw) between the predicted orientation
and the ground truth.

\subsubsection{Absolute Yaw Error}
The equations below show how to calculate the absolute yaw error (AYE).

\begin{equation}
    \psi = \arctan2\!\bigl(2(wz + xy),\; 1 - 2(y^2 + z^2)\bigr)
\end{equation}

\begin{equation}
    \epsilon_{\text{yaw}} =
    \begin{cases}
        |\psi_{\text{pred}} - \psi_{\text{gt}}| & \text{if } |\psi_{\text{pred}} - \psi_{\text{gt}}| \leq 180^{\circ} \\
        360^{\circ} - |\psi_{\text{pred}} - \psi_{\text{gt}}| & \text{otherwise}
    \end{cases}
\end{equation}

The yaw is extracted from the quaternion rotation before the use of the equation above. Further, the AYE is wrapped over [0, 180], such that errors near $\pm180^{\circ}$ are not artificially inflated.

The reason as to why the AYE is an interesting metric to track is because of \fixme{mag section}. In this section it is mentioned that only the magnetometer can provide
a heading reference and thus with corrupted magnetometer measurements this drift in yaw becomes unbounded. By specifically looking at AYE, we can see how much the proposed solution is
able to correct yaw especially due to the test set conditions in \fixme{train/test/val section}. 

\subsection{Test Pipeline}

There a two different pipelines tested within this project. First, is dead reckoning (gyroscope integration only), this pipeline tests how the TCN performs by only providing $\omega$ corrections
which is applied additively to the raw gyroscope measurement. This test shows how the TCN performs solely as a gyroscope corrector. Second, is the full proposed solution with the integration of 
the MUKF. This tests shows the performance of the corrected gyroscope measurement in tandem with corrections applied from the accelerometer and magnetometer by the update step in the MUKF.

The rest of the results section will cover these two pipelines and report the per file statistics and the total aggregated statistics over the entire test set. 